\شروع‌مسئله[۱۰] \متن‌سیاه{جمع صفر}.
$2k$
آرایه داریم که هر کدام شامل $n$ عدد صحیح است. الگوریتمی طراحی کنید که از هر آرایه یک عدد (یعنی در کل
$2k$
عدد) انتخاب کند که مجموع این اعداد برابر با صفر شود و یا اعلام کند که چنین انتخابی ممکن نیست. زمان اجرای الگوریتم شما باید حداکثر $O(n^k)$ باشد. البته به الگوریتم با زمان
$O(k\cdot n^k)$
نیز حداکثر ۷ نمره خواهد گرفت.

\vspace{0.5cm}

\شروع‌مسئله[۱۰] \متن‌سیاه{قرمز آبی}.
یک گراف جهت‌دار داریم که هر کدام از یال‌های آن با یکی از رنگ‌های قرمز و یا آبی رنگ شده است. ما به دنبال پیدا کردن یک دور جهت‌دار هستیم که تعداد یال‌های قرمز آن اکیداً بیشتر از تعداد یال‌های آبی آن باشد. در حالتی که چندین دور وجود داشته باشد، پیدا کردن یکی از این آن‌ها کفایت می‌کند و در صورتی که چنین دوری وجود نداشته، می‌بایست این موضوع را تشخیص دهیم.

می‌دانیم در صورتی که به‌طور مستقیم از الگوریتم بلمن‌فورد استفاده کنیم، چنین دوری را می‌توان در زمان $O(ne)$ پیدا کرد که $n$ تعداد رأس‌های و $e$ تعداد یال‌های گراف است. 
اگر تعداد یال‌های گراف نیز $O(n)$ باشد، این زمان خود به خود به $O(n^2)$ بهبود پیدا می‌کند.
اما ما به دنبال الگوریتم سریع‌تری هستیم چون گراف ما خاصیت دیگری نیز دارد.

در این گراف به رأس‌هایی که جمع درجه ورودی و خروجی آن بیش از ۲ باشد، رأس شلوغ گفته می‌شود. تعداد رأس‌های شلوغ در این گراف برابر با $k$ است که کمتر از $n$ است (مثلاً $k=O(\sqrt{n})$). همچنین می‌دانیم که تعداد یال‌های گراف $O(n)$ است. الگوریتمی طراحی کنید که در زمان $O(n+k^3)$ این مسئله را حل کند. 

\صفحه‌جدید

\شروع‌مسئله[۱۰] \متن‌سیاه{استقلال پرسپولیس}.
در دوران قبل از کرونا، ملانصرالدین برای دیدن بازی شهرآورد وارد تهران می‌شود و از فرودگاه تا استادیوم را می‌خواهد پیاده‌روی کند. خیابان‌های شهر به سه دسته تقسیم شده‌اند: خیابان‌هایی که طرفداران استقلال در آن غالب هستند، خیابان‌هایی که طرفداران پرسپولیس در آن غالب هستند و خیابان‌هایی که طرفداران هیچ‌کدام از دو تیم در آن غالب نیستند. ملانصرالدین وقتی می‌خواهد از خیابانی عبور کند، در صورتی که طرفداران استقلال در آن خیابان غالب باشند، می‌خواهد پیراهن آبی به تن داشته باشد و در صورتی که طرفداران پرسپولیس در آن غالب باشند، می‌خواهد پیراهن قرمز به تن داشته باشد. در حالت سوم رنگ پیراهنش می‌تواند آبی یا قرمز باشد. ملانصرالدین از ترس این که لو برود که پیراهنش را عوض می‌کند، اگر پیراهن جدیدی به تن کند، پیراهن قبلی‌اش را دور می‌اندازد! او در ابتدا نیز هیچ پیراهنی به همراه ندارد و با لباس سنتی وارد فرودگاه می‌شود.
در فرودگاه و تعدادی از تقاطع‌های شهر فروشگاه‌های لباس وجود دارد که ملانصرالدین می‌تواند پیراهن قرمز یا آبی از آن خریداری کند.

نقشه شهر به این صورت است که از $m$ خیابان دوطرفه تشکیل شده که هر خیابان بین دو تقاطع قرار دارد. تعداد تقاطع‌ها نیز $n$ است و به هر تقاطع حداکثر ۴ خیابان متصل است. فرودگاه و استادیوم را نیز می‌توان تقاطع به حساب آورد.
شما با داشتن تمامی اطلاعات شامل نقشه شهر، لیست تقاطع‌هایی که در آن‌ها فروشگاه لباس وجود دارد و طول هر خیابان، باید الگوریتمی با زمان $O(n\log n)$ طراحی کنید که کوتاهترین مسیر از فرودگاه به استادیوم را به ملانصرالدین پیشنهاد دهد. همچنین فرض کنید  خرید پیراهن جدید از فروشگاه لباس و تعویض پیراهن معادل رد شدن از یک خیابان ۱۰۰ متری زمان‌بر است که در محاسبه طول یک مسیر درنظر گرفته می‌شود. همچنین فرض کنید ملانصرالدین برای خرید پیراهن‌های متعدد محدودیت مالی ندارد.

\پاسخ{
	
}


